\documentclass[11pt]{article}
\usepackage{mathblatt}
% gesetzt von werkzeuge/setzer.py; Lösungen nach bau/bauregeln.md „Lösungen“.
\newcommand{\kennung}{DXZ}
\begin{document}
\pfheftstil
\fancyfoot[C]{{\footnotesize\color{mbgrau}\kennung\hfill\thepage}}
\pfheftkopf{Dezimalzahlen}{Klasse 6 \textperiodcentered{} Lösungen}

\begin{pfloesung}{}
\lz{1b)}{a) $0{,}3$ \quad b) $0{,}9$ \quad c) $0{,}47$ \quad d) $0{,}81$}{a) $0{,}3$; b) $0{,}9$; c) $0{,}47$; d) $0{,}81$}{}
\end{pfloesung}
\begin{pfloesung}{}
\lz{2.}{a) $0{,}06$ \quad b) $0{,}006$ \quad c) $0{,}052$ \quad d) $0{,}09$}{a) $0{,}06$; b) $0{,}006$; c) $0{,}052$; d) $0{,}09$}{}
\end{pfloesung}
\begin{pfloesung}{}
\lz{3.}{$4{,}058 = \frac{4058}{1000}$}{}{}
\end{pfloesung}
\begin{pfloesung}{}
\lz{4.}{a) $\frac{7}{10}$ \quad b) $\frac{13}{100}$ \quad c) $\frac{8}{100}$ \quad d) $\frac{2019}{1000}$}{a) $\frac{7}{10}$; b) $\frac{13}{100}$; c) $\frac{8}{100}$; d) $\frac{2019}{1000}$}{}
\end{pfloesung}
\begin{pfloesung}{}
\lz{5.}{$12{,}07$\,s; Ben}{Ben war schneller: Ali braucht 12{,}7\,s, und 7 Zehntel sind mehr als 7 Hundertstel.}{}
\end{pfloesung}
\begin{pfloesung}{}
\lz{6b)}{$A = 0{,}4$, \ $B = 0{,}9$, \ $C = 1{,}3$}{Ein Schritt ist $0{,}1$: $A = 0{,}4$, $B = 0{,}9$, $C = 1{,}3$}{}
\end{pfloesung}
\begin{pfloesung}{}
\lz{7.}{3 hinter $0{,}4$; \ 1 vor $0{,}5$; \ 6 hinter $0{,}5$}{$0{,}43$: 3 Striche hinter $0{,}4$; $0{,}49$: 1 Strich vor $0{,}5$; $0{,}56$: 6 Striche hinter $0{,}5$}{}
\end{pfloesung}
\begin{pfloesung}{}
\lz{8.}{a) $3{,}4$ | $3{,}5$ \quad b) $0{,}7$ | $0{,}8$ \quad c) $5{,}0$ | $5{,}1$ \quad d) $1{,}9$ | $2{,}0$}{a) $3{,}4$ und $3{,}5$; b) $0{,}7$ und $0{,}8$; c) $5{,}0$ und $5{,}1$; d) $1{,}9$ und $2{,}0$}{}
\end{pfloesung}
\begin{pfloesung}{}
\lz{9.}{a) $2{,}9; 3{,}0; 3{,}1$ \quad b) $1{,}0; 1{,}2; 1{,}4$ \quad c) $4{,}99; 5{,}00; 5{,}01$}{a) $2{,}9; 3{,}0; 3{,}1$; b) $1{,}0; 1{,}2; 1{,}4$; c) $4{,}99; 5{,}00; 5{,}01$}{}
\end{pfloesung}
\begin{pfloesung}{}
\lz{10.}{$3{,}4$\,m und $3{,}5$\,m; näher an $3{,}5$\,m}{Zwischen $3{,}4$\,m und $3{,}5$\,m (10\,cm $= 0{,}1$\,m); näher an $3{,}5$\,m: bis dahin 4\,cm, bis $3{,}4$\,m sind es 6\,cm.}{}
\end{pfloesung}
\begin{pfloesung}{}
\lz{11b)}{a) $0{,}5$ \quad b) $0{,}6$ \quad c) $0{,}8$ \quad d) $1{,}5$}{a) $\frac{5}{10} = 0{,}5$; b) $\frac{6}{10} = 0{,}6$; c) $\frac{8}{10} = 0{,}8$; d) $\frac{15}{10} = 1{,}5$}{}
\end{pfloesung}
\begin{pfloesung}{}
\lz{12.}{a) $0{,}75$ \quad b) $0{,}45$ \quad c) $0{,}24$ \quad d) $0{,}14$}{a) $\frac{75}{100} = 0{,}75$; b) $\frac{45}{100} = 0{,}45$; c) $\frac{24}{100} = 0{,}24$; d) $\frac{14}{100} = 0{,}14$}{}
\end{pfloesung}
\begin{pfloesung}{}
\lz{13.}{a) $\frac{2}{5}$ \quad b) $\frac{1}{4}$ \quad c) $\frac{13}{20}$ \quad d) $\frac{6}{5}$}{a) $\frac{4}{10} = \frac{2}{5}$; b) $\frac{25}{100} = \frac{1}{4}$; c) $\frac{65}{100} = \frac{13}{20}$; d) $\frac{12}{10} = \frac{6}{5}$}{}
\end{pfloesung}
\begin{pfloesung}{}
\lz{14.}{$\frac{3}{5} = 0{,}6$, \ $\frac{7}{20} = 0{,}35$, \ $\frac{1}{4} = 0{,}25$, \ $\frac{2}{25} = 0{,}08$}{$\frac{3}{5} = 0{,}6$; $\frac{7}{20} = 0{,}35$; $\frac{1}{4} = 0{,}25$; $\frac{2}{25} = 0{,}08$}{}
\end{pfloesung}
\begin{pfloesung}{}
\lz{15.}{50\,g fehlen}{$\frac{3}{4}$\,kg $= \frac{75}{100}$\,kg $= 0{,}75$\,kg $= 750$\,g; $0{,}7$\,kg $= 700$\,g; es fehlen $750 - 700 = 50$\,g.}{}
\end{pfloesung}
\begin{pfloesung}{}
\lz{16b)}{a) $0{,}25$ \quad b) $0{,}75$ \quad c) $0{,}125$ \quad d) $0{,}875$}{a) $1 : 4 = 0{,}25$; b) $3 : 4 = 0{,}75$; c) $1 : 8 = 0{,}125$; d) $7 : 8 = 0{,}875$}{}
\end{pfloesung}
\begin{pfloesung}{}
\lz{17.}{a) $2{,}25$ \quad b) $1{,}375$ \quad c) $2{,}625$}{a) $9 : 4 = 2{,}25$; b) $11 : 8 = 1{,}375$; c) $21 : 8 = 2{,}625$}{}
\end{pfloesung}
\begin{pfloesung}{}
\lz{18.}{a) $0{,}\overline{3}$ \quad b) $0{,}\overline{6}$ \quad c) $0{,}\overline{5}$ \quad d) $0{,}1\overline{6}$}{a) $0{,}\overline{3}$; b) $0{,}\overline{6}$; c) $0{,}\overline{5}$; d) $0{,}1\overline{6}$}{}
\end{pfloesung}
\begin{pfloesung}{}
\lz{19.}{a) ja \quad b) ja \quad c) nein \quad d) ja}{a) ja, $0{,}375$; b) ja, $0{,}35$; c) nein, $5 : 6 = 0{,}8333\ldots = 0{,}8\overline{3}$; d) ja, $0{,}36$}{}
\end{pfloesung}
\begin{pfloesung}{}
\lz{20.}{Ja, $62{,}5$\,cm je Stück}{Ja; ein Stück ist $\frac{5}{8}$\,m, $5 : 8 = 0{,}625$, also $0{,}625$\,m $= 62{,}5$\,cm. Das sind $2{,}5$\,cm mehr als 60\,cm.}{}
\end{pfloesung}
\begin{pfloesung}{}
\lz{21.}{a) $0{,}9$ \quad b) $0{,}09$ \quad c) $0{,}045$}{a) $0{,}9$; b) $0{,}09$; c) $0{,}045$}{}
\end{pfloesung}
\begin{pfloesung}{}
\lz{22.}{$A = 1{,}26$, \ $B = 1{,}35$; \ $1{,}2$ | $1{,}3$}{Ein Schritt ist $0{,}01$: $A = 1{,}26$, $B = 1{,}35$; $A$ liegt zwischen $1{,}2$ und $1{,}3$}{}
\end{pfloesung}
\begin{pfloesung}{}
\lz{23.}{a) $0{,}6$ \quad b) $0{,}85$ \quad c) $\frac{9}{20}$}{a) $\frac{6}{10} = 0{,}6$; b) $\frac{85}{100} = 0{,}85$; c) $\frac{45}{100} = \frac{9}{20}$}{}
\end{pfloesung}
\begin{pfloesung}{}
\lz{24.}{a) $0{,}875$ \quad b) $0{,}\overline{6}$}{a) $7 : 8 = 0{,}875$; b) $2 : 3 = 0{,}666\ldots = 0{,}\overline{6}$}{}
\end{pfloesung}
\begin{pfloesung}{}
\lz{25.}{$\frac{7}{20} = 0{,}35$}{$\frac{7}{20} = \frac{35}{100}$; richtig wäre $\frac{1}{5} = 0{,}2$, $\frac{3}{100} = 0{,}03$, $\frac{1}{3} = 0{,}\overline{3}$.}{}
\end{pfloesung}
\begin{pfloesung}{}
\lz{26.}{25\,g zu viel}{Zu viel; $\frac{3}{8}$\,kg: $3 : 8 = 0{,}375$, also 375\,g; $0{,}4$\,kg $= 400$\,g; $400 - 375 = 25$\,g zu viel.}{}
\end{pfloesung}
\end{document}
